\section{Specialization to Reidemeister-III unlocking}
\label{sec:r3}

We now verify the abstract hypotheses for the mutable dart representation used
by the reviewed ProveIt code.  The purpose is not to reprove Reidemeister's
theorem; it is to account for the local search tree.

\subsection{Fixed darts and the face permutation}

An $N$-crossing diagram has darts
\[
  V=\{0,1,\ldots,4N-1\}.
\]
Darts $4i,\ldots,4i+3$ belong to crossing $i$.  Let $\alpha$ be the fixed-point
free involution pairing the two darts of each diagram edge.  Let $\rho$ rotate a
dart to the next slot at its crossing.  A face walk is generated by
\[
  \varphi=\rho\circ\alpha.
\]
The implementation writes this map as \code{nxt}.  A triangular face is a
three-cycle of $\varphi$ through three distinct crossings.  The over/under slot
parities determine whether that triangle is an admissible RIII configuration.

The RIII update keeps crossing identities and slots fixed and modifies only the
edge involution.  This is particularly convenient for the causal analysis:
site names do not move even though adjacency does.

\subsection{A conservative RIII footprint}

Let $T$ be a legal RIII triangle with crossing set $C(T)$ of size three.  Define
\[
  D(T)=\{4c+j:c\in C(T),\ 0\leq j<4\}
\]
and the conservative footprint
\[
  \Phi(T)=D(T)\cup\alpha(D(T)).
\]
Thus $|\Phi(T)|\leq24$.

\begin{lemma}[Footprint containment]
\label{lem:r3-footprint}
The reviewed \code{apply\_r3} routine changes $\alpha$ only on darts in
$\Phi(T)$.  Its inverse has the same footprint set.
\end{lemma}

\begin{proof}
The routine names the six outer slots of the three participating strands and
the corresponding inner slots.  Every new link has one endpoint at a dart of
$D(T)$; its other endpoint is either another dart of $D(T)$ or the previous
$\alpha$-partner of a dart in $D(T)$.  Hence every written array entry lies in
$D(T)\cup\alpha(D(T))$.  Moreover, each outside dart in
$\alpha(D(T))\setminus D(T)$ remains paired to some dart of $D(T)$ after the
move, while no new outside dart is introduced.  Therefore the post-move set
$D(T)\cup\alpha'(D(T))$ equals the pre-move footprint.  RIII is an involution,
so the reverse action has exactly this same footprint.
\end{proof}

For RI/RII terminal tests, use the analogous conservative footprint consisting
of all darts at the one or two participating crossings together with their edge
partners.  Its size is bounded by 16.

\subsection{Disjoint footprints commute}

\begin{proposition}[Strong locality for the dart rules]
\label{prop:r3-locality}
RIII actions and RI/RII terminal tests in the fixed-dart implementation satisfy
strong disjoint commutation when the conservative footprints above are used.
\end{proposition}

\begin{proof}
An RIII action writes only entries of $\alpha$ in its footprint by
Lemma~\ref{lem:r3-footprint}.  Suppose a second RIII action is observed after
the first and their residual footprints are disjoint.  The twelve darts owned
by the second crossing triple lie outside the first footprint.  Since the first
move changes $\alpha$ only inside its own footprint, the partners of those
twelve darts are also unchanged; hence the second conservative footprint, its
three-cycle face test, and its slot-parity test were already identical before
the first move.  The same argument in the other direction proves the two-sided
legality diamond.  The two updates write disjoint array entries, so they commute
entry by entry and their footprints remain fixed under the swap.

The same argument applies to RI/RII legality: monogon and bigon face walks read
only the darts at their participating crossings and their partners.  A
disjoint RIII update neither creates nor destroys such a local face.  Equivalently,
any RI/RII move newly exposed by an RIII update has a conservative footprint
intersecting that RIII footprint.  Finally, RIII inversion is its own local
inverse with the same footprint.
\end{proof}

\begin{remark}
Two RIII moves with disjoint \emph{crossing triples} need not have disjoint
conservative footprints, because their local disks can touch the same outside
edge.  This is why the theorem is formulated with dart footprints rather than
only crossing labels.
\end{remark}

\subsection{Global and local incidence}

For a connected spherical $4$-regular projection graph with $N$ vertices,
$E=2N$ and Euler's formula gives $F=N+2$ faces.  A legal RIII action is supported
by a triangular face, so the total number of legal RIII locations is at most
$N+2$.  This gives the global incidence bound with an absolute $A$.

For the local bound, fix one dart $d$.  If $d\in\Phi(T)$, then either a crossing
of $T$ owns $d$, or a crossing of $T$ owns $\alpha(d)$.  Each crossing has four
face corners, so it is incident to at most four triangular faces.  Consequently
at most eight triangular faces have footprints containing a given dart.  By a
union bound,
\[
  |\{T:\Phi(T)\cap S\neq\varnothing\}|\leq 8|S|.
\]
Thus one may take $f=24$ and $B=8$, giving $Bf=192$.  Whenever an RIII
location exists, $N\geq3$ and $N+2\leq2N$, so a canonical enumeration with one
representative per triangular face may take $A=2$ and hence $C=384$ in
Theorem~\ref{thm:enumeration}.  An implementation that emits orientations or
duplicate representatives is covered by some larger absolute $C$; no numerical
optimality is claimed, and the asymptotic regimes are unchanged.

\begin{corollary}[RIII search bound]
\label{cor:r3-bound}
On an $N$-crossing fixed-dart state, a first-unlocking RIII trace of length at most $k$ and birth number at most
$b$, when one exists, can be found by exhaustive birth-front search in
\[
  N^{b+O(1)}(Ck)^k
\]
time and polynomial memory for an absolute constant $C$.
\end{corollary}

\subsection{Why accumulated footprints are the right implementation object}

The existing recursive search computes a neighborhood from the triangle just
applied and recursively scans that neighborhood.  This is a last-touch chain.
For birth number one, the complete normal-form search must instead retain the
union of all previous RIII footprints.  A candidate is local when its current
footprint intersects that union.  The active set contains at most $24k$ darts,
so scanning crossings owning active darts and their current partners generates
only $O(k)$ local triangle candidates, up to duplicates.

For $b>1$, all births are chosen first.  They are required to have pairwise
disjoint footprints and are ordered canonically.  After the first connected
move, no further birth is permitted.  This exactly implements
Theorem~\ref{thm:birth-front}.

\subsection{Interaction with the face shortcut}

The current implementation's \code{r3\_can\_help} test is used only for the
last allowed RIII move.  It checks a necessary local condition for that move to
create an RI/RII face.  Retaining that placement is important.  Applying the
shortcut to every intermediate move would be unsound, because an intermediate
RIII need not immediately expose a reduction; it may only prepare a later
interaction.

\subsection{State restoration}

Every unsuccessful branch must restore the exact involution.  The reviewed code
undoes an RIII move by locating the inverse triangle and applying the same
involution routine.  The adapter preserves this contract.  A production
integration should additionally test every budget-exhaustion boundary, because
an early return from a recursive branch is a common place for mutable-state
leaks.
